\documentclass{article}
\usepackage{amsmath}
\begin{document}
\title{Waiting for the first qualified camera frame consumes motion slack, not just notice time}
\maketitle

\paragraph{Interface.}
Retain the far site $x=3.9$, $q=\sigma-.15$, $\sigma=4/h$,
$h=6.37$, $\alpha=10$, and the unit speed bound. At the end of the
nominal dwell the robot is at rest. If the next valid frame is
captured only at $q+\epsilon$, the robot must remain stationary
until then. A qualified capture-completion mark is observable at
release; the true world bit is separately usable after an additional
processing delay $L$. This is a two-cut specified detector, not an
assumption of calibrated field timing or a proof that a native
camera call will expose at any requested time.

\paragraph{Exact charged slack.}
Write $\delta=7-h-\sigma\simeq .002056515$. Both worlds are
feasible exactly when
\[
  0\le\epsilon\le\delta,\quad
 0\le L \le .1 + \sqrt{(\delta-\epsilon)/\alpha}.
\]
For $\epsilon>\delta$, even a known far world is impossible:
the pointwise highest velocity starting from rest at $q+\epsilon$
gives $x(2)\le4+2-\sigma-\epsilon<h-1$. This argument allows
any motion consistent with the actuator; it is not a failure
of a specific planner.

For the allowed waits put $z=\sqrt{\alpha(\delta-\epsilon)}$ and
$v_*=1-z$. The exact last branch point is
$r_*=q+\epsilon+.1+z/\alpha$. From capture, accelerate to unit
speed, then brake to $v_*$ at $r_*$. In the far branch resume
positive acceleration then coast, losing exactly $z^2/\alpha$
of range compared to the maximum. This reaches $h-1$ at time 2,
and $x(\sigma)\le4$ automatically from the maximum-speed envelope
starting at the station. The far coverage follows as in the parent.
For late, continue braking from unit speed to rest exactly at $4$
at time $q+\epsilon+.2>2/3$, then hold. An earlier bit can be
held until the same calibrated branch point.

For necessity at a cut $r$ slightly beyond $r_*$, ending the common
segment at $v_*$ allows a unit-speed plateau: integrate the minimum
of the acceleration cones and the speed ceiling. The combined
common-plus-far maximum with this $v_*$ is exactly $h-1$ and does
not increase with $r$. It is strictly smaller for $v<v_*$, since
the far displacement is strictly increasing on that range and the
common bound is nondecreasing. Thus $v\ge v_*$. The late stopping
peak is after its seam, since $r_*+v_*/\alpha=q+\epsilon+.2>2/3$.
The common position must fit below $4-v^2/(2\alpha)$ to stop safely
and above $h-1-(2-r)+(1-v)^2/(2\alpha)$ to attain the far target.
Together those require $r\le7-h-[v^2+(1-v)^2]/(2\alpha)\le r_*$
because $v_*>1/2$. The usual ignore-early-bit argument rules out
any even later feasible receipt. This proves the exact relative
processing bound stated above.

\paragraph{Native rate gate and a design comparison.}
The pinned CoSTAR SDF specifies 20 Hz. Execute the pinned
\texttt{Sensor::Update(now,force)} body with a stand-in nominal time
and a counted callback. The regular gate skips before
\texttt{nextUpdateTime} and adds a period to that planned time
when admitted; it does not reset the phase to the actual call
time. The forced arm bypasses the regular gate in the fixture
and does not advance that planned time. Calls of the underlying
renderer, pose read and actuation rights are separate; production
packages are not asserted to expose any authorized external force
operation. A known aligned regular phase is another adequate
comparison, not a capability gap versus a forced call.

A nominal unknown phase over 50 ms cannot guarantee this far task.
Even zero processing leaves only $\delta$ of stationary waiting.
At modeled processing .105, the allowable wait is
$\delta-.00025\simeq .001806515$. These are not empirical phase
frequencies. For 1,000 audited phases at 50 microsecond resolution,
the native gate admits 37 into that physical service envelope at
.105. Forcing at the pinned eligible capture would pass the timing
test in that arm; it does not fix processing .125. The independent
motion check uses exact rational paths between the analytical
boundaries. Safety, frame validity and independent mission
verification remain open.
\end{document}
